% !TEX program = xelatex
% SG-Lean 思路汇报
% 编译：xelatex sg-lean-report.tex  （连续两遍以稳定交叉引用）
\documentclass[UTF8,a4paper,11pt,fontset=windows]{ctexart}

\usepackage{geometry}
\geometry{left=2.4cm,right=2.4cm,top=2.6cm,bottom=2.4cm}
\usepackage{amsmath,amssymb}
\usepackage{booktabs}
\usepackage{tabularx}
\usepackage{array}
\usepackage{enumitem}
\usepackage{xcolor}
\usepackage{tikz}
\usetikzlibrary{positioning,arrows.meta,fit,backgrounds,calc}
\usepackage{tcolorbox}
\tcbuselibrary{skins,breakable}
\usepackage{fancyhdr}
\usepackage{url}
\usepackage{adjustbox}
\usepackage[hidelinks]{hyperref}

\definecolor{sgblue}{RGB}{31,86,153}
\definecolor{sggreen}{RGB}{38,120,86}
\definecolor{sggray}{RGB}{90,96,106}

% 圈码 ①-⑳ 属于 Enclosed Alphanumerics，默认走拉丁字体会缺字，
% 显式划入 CJK 字符类，交给中文字体渲染。
\xeCJKDeclareCharClass{CJK}{"2460 -> "24FF}

\ctexset{
  section/format = \Large\bfseries\color{sgblue},
  subsection/format = \large\bfseries\color{sgblue!85!black},
}

\pagestyle{fancy}
\fancyhf{}
\fancyhead[L]{\small\color{sggray}SG-Lean 思路汇报}
\fancyhead[R]{\small\color{sggray}\leftmark}
\fancyfoot[C]{\small\thepage}
\renewcommand{\headrulewidth}{0.4pt}

\newcolumntype{Y}{>{\raggedright\arraybackslash}X}
\newcommand{\lp}[1]{\path{#1}}

\tikzset{
  flowbox/.style = {draw, rounded corners=2pt, align=center, inner sep=4pt,
                    minimum height=7mm, fill=blue!5, font=\small},
  sidebox/.style = {draw, rounded corners=2pt, align=center, inner sep=4pt,
                    fill=orange!10, font=\footnotesize},
  offbox/.style  = {draw, rounded corners=2pt, align=center, inner sep=4pt,
                    fill=green!6, font=\small, minimum height=7mm},
  arr/.style     = {-{Latex[length=2mm]}, thick, draw=sggray},
  dasharr/.style = {-{Latex[length=2mm]}, thick, dashed, draw=sggray},
}

\newtcolorbox{keybox}[1][]{
  colback=blue!4, colframe=sgblue!70, boxrule=0.6pt, arc=2pt,
  left=6pt, right=6pt, top=5pt, bottom=5pt, breakable, #1}

\title{\bfseries SG-Lean 思路汇报\\[2mm]
       \large 面向 Lean 4 的形式化数学证明器}
\author{}
\date{}

\begin{document}
\maketitle
\thispagestyle{fancy}

\begin{keybox}
\textbf{文档说明}：本文件说明项目的目标、总体思路、执行流程与工具模块，
是当前思路的权威口径。仓库中若干历史文档（\texttt{README.md}、
\texttt{docs/framework.md} 等）描述的是代码实现的演变过程，若与本文件冲突，以本文件为准。
\end{keybox}

\vspace{2mm}
\tableofcontents
\vspace{4mm}

%======================================================================
\section{目标与指标}

\subsection{项目目标}

本项目的目标是构建一个 \textbf{Lean 4 形式化数学证明器}：输入一条数学命题的
Lean 形式化陈述，输出一段能够通过 Lean 内核检验的证明脚本。

\begin{table}[htbp]
\centering
\small
\begin{tabularx}{\textwidth}{@{}p{2.2cm}Y@{}}
\toprule
\textbf{输入} & 一条 Lean 命题。可以是 \texttt{theorem} 或 \texttt{example} 语句
（其证明体为 \texttt{sorry}），从语句头部抽取绑定变量转换为全称量化的闭式命题；
也可以是等价的闭式命题（只引用全局常量）。\\
\midrule
\textbf{输出} & 一段 tactic 证明脚本，外加验证记录（内核实测通过、耗时、来源后端等）。
\\
\midrule
\textbf{保证} & 输出的每一篇证明都经过 Lean 内核终检。未通过验证的候选证明
\textbf{不会}出现在输出中。\\
\bottomrule
\end{tabularx}
\end{table}

\subsection{指标与优先级}

项目的全部工程取舍，按下列顺序排列优先级：

\begin{table}[htbp]
\centering
\small
\begin{tabularx}{\textwidth}{@{}p{1.0cm}p{2.8cm}Y@{}}
\toprule
\textbf{优先级} & \textbf{指标} & \textbf{定义与说明} \\
\midrule
1 & 正确率（主指标） & 在固定评测集上报 \textbf{pass@1} 与 \textbf{pass@k}，
\textbf{目标级}统计：一道题只要有至少一篇候选通过内核检验，即算该题被解出。 \\
\midrule
2 & 模型调用成本 & token 消耗量与调用次数。任何新增杠杆都必须报出它的边际成本。 \\
\midrule
3 & 时间成本 & 墙钟时间，含 Lean 侧验证的等待（Lean 批处理的固定开销较大，需单独核算）。
\\
\bottomrule
\end{tabularx}
\end{table}

这一顺序意味着：当两条技术路线在正确率上无法区分时，选成本低的那条；
当一条路线要花更多调用换取少量正确率提升时，需要有实测数据支撑才会被采纳。

\subsection{范围与边界}

\begin{itemize}[leftmargin=1.6em,itemsep=2pt]
  \item \textbf{不做模型训练}：不训练、不微调基座模型。证明能力的上限由求解模型决定，
        本项目贡献的是模型之上的脚手架，以及用 Lean 保证的输出正确性。
  \item \textbf{不做自动形式化}：不处理自然语言题干到 Lean 命题的转换，输入即为 Lean 命题。
  \item \textbf{不承诺通用性}：只承诺在选定评测基准与 Mathlib 覆盖范围之内尽可能高的正确率；
        超出范围的命题能力仅由基座模型决定。
  \item \textbf{搜索暂为储备}：证明搜索（MCTS 一类）作为可选后端保留，当前不作为主线。
\end{itemize}

%======================================================================
\section{思路框架与执行流程}

\subsection{总体框架}

整体上，本项目由\textbf{两条时间尺度}构成，二者通过\textbf{引理库}这一个接口相连：

\begin{itemize}[leftmargin=1.6em,itemsep=2pt]
  \item \textbf{在线求解}（秒级至分钟级）：接收一道题，输出一篇通过验证的证明。
  \item \textbf{离线资产生产}（小时级）：在一批课程集上运行闭环，把重复出现的能力缺口
        转化为已通过验证的引理，沉淀为库。
\end{itemize}

\begin{figure}[htbp]
\centering
\begin{adjustbox}{max width=\textwidth}
\begin{tikzpicture}[font=\small,node distance=5mm]
  \node[flowbox,fill=blue!8] (q) {命题};
  \node[flowbox,fill=blue!8,right=6mm of q] (r) {分层检索};
  \node[flowbox,fill=blue!8,right=6mm of r] (s) {证明后端\\（可插拔）};
  \node[flowbox,fill=blue!8,right=6mm of s] (v) {内核验证};
  \node[flowbox,fill=blue!8,right=6mm of v] (p) {证明};
  \draw[arr] (q)--(r); \draw[arr] (r)--(s); \draw[arr] (s)--(v); \draw[arr] (v)--(p);
  \node[draw,dashed,rounded corners=4pt,inner sep=4mm,fit=(q)(p)] (onl) {};
  \node[above=0.5mm of onl,font=\small\bfseries,text=sgblue]
        {在线求解：一道题（秒级至分钟级）};

  \node[offbox,below=30mm of q] (c) {课程集 $C$};
  \node[offbox,right=6mm of c] (d) {需求挖掘};
  \node[offbox,right=6mm of d] (g) {引理生成};
  \node[offbox,right=6mm of g] (h) {形式化判据};
  \node[offbox,right=6mm of h] (l) {引理库};
  \draw[arr] (c)--(d); \draw[arr] (d)--(g); \draw[arr] (g)--(h); \draw[arr] (h)--(l);
  \node[draw,dashed,rounded corners=4pt,inner sep=4mm,fit=(c)(l)] (off) {};
  \node[below=0.5mm of off,font=\small\bfseries,text=sggreen]
        {离线资产生产：一批课程集（小时级）};

  \draw[arr] (off.north) -- (onl.south)
        node[midway,right=1mm,font=\footnotesize,text=sggray]{引理库注入（资产）};
\end{tikzpicture}
\end{adjustbox}
\caption{总体框架：在线求解与离线资产生产，通过引理库相连。}
\end{figure}

\textbf{三条不变量}（框架的骨架，不做妥协）：

\begin{keybox}
\begin{enumerate}[leftmargin=1.6em,itemsep=3pt]
  \item \textbf{能力来自基座模型}。我们不训练大模型，因此不承诺通用性，
        只承诺在基准与 Mathlib 覆盖范围内尽可能高的正确率。
  \item \textbf{增益来自可寻址的资产}，而不是权重。资产具备两个性质：\textbf{可加}
        （新增一个领域的引理不会破坏已有领域的引理）与\textbf{可寻址}
        （解决一道题时只需要其中相关的少数几条）。这是应对未知题型的结构性优势。
  \item \textbf{正确性由 Lean 内核保证}，不由模型保证。任何来源的候选证明，
        都必须经过同一道内核闸门。
\end{enumerate}
\end{keybox}

\subsection{与 SGS 的关系}

本项目的思路源于 SGS（\emph{Scaling Self-Play with Self-Guidance}）的三角色自博弈框架：
求解者（Solver）、出题者（Conjecturer）与评审者（Guide）。
我们保留其\textbf{结构、循环形状与防退化思想}，替换其中两样东西：

\begin{table}[htbp]
\centering
\small
\begin{tabularx}{\textwidth}{@{}p{4.6cm}Y@{}}
\toprule
\textbf{SGS（训练期）} & \textbf{本项目（推理期）} \\
\midrule
Guide 由大模型对合成题打分 & 替换为\textbf{形式化判据}
$G = \text{硬门} \times \text{软分}$，全部由 Lean 计算，确定性、不消耗 token \\
\midrule
用强化学习更新求解者与出题者的权重 & 替换为\textbf{引理库的更新}，
并以提示词注入的方式在下一轮被使用 \\
\midrule
出题者条件化在未解目标上 & 保留，并扩展输入：未解目标 $+$ 需求签名 $+$ 库中已入库引理范例 \\
\midrule
奖励为 $R_{\text{solve}} \cdot R_{\text{guide}}$ & 保留乘法结构：可证性以硬门形式保留，
质量以软分（压缩收益与覆盖收益）体现 \\
\bottomrule
\end{tabularx}
\end{table}

\textbf{两点必须如实说明}：

\begin{enumerate}[leftmargin=1.6em,itemsep=3pt]
  \item SGS 的循环\textbf{依靠梯度闭合}。本项目不训练权重，因此本项目的形态对应于 SGS
        消融实验中的“冻结出题者”一档。SGS 论文的原话是其表现优于纯强化学习基线，
        但劣于完整方法。我们以“资产累积”承担原本由权重更新承担的角色，
        这一替代是否充分，正是本项目要回答的问题。
  \item SGS 的评测口径是在\textbf{训练题}上统计累计解出率；
        本项目采用留出集协议（课程集建库、开发集调参、测试集冻结后只跑一次），
        在评测严谨性上比 SGS 更严格。
\end{enumerate}

\subsection{执行流程一：在线求解一道题}

\begin{figure}[htbp]
\centering
\begin{adjustbox}{max width=\textwidth}
\begin{tikzpicture}[font=\small,node distance=6mm]
  \node[flowbox,text width=62mm] (in) {输入：一条 Lean 命题（闭式陈述）};
  \node[flowbox,text width=62mm,below=of in] (g)
        {① 门检：将命题 elaborate 为 \texttt{Prop}};
  \node[flowbox,text width=62mm,below=of g] (t)
        {② 廉价 tactic 兜底：\texttt{decide} / \texttt{simp} / \texttt{norm\_num} / \texttt{ring} 等};
  \node[flowbox,text width=62mm,below=of t] (r)
        {③ 分层检索：引理库（精排层）$+$ Mathlib（兜底层）};
  \node[flowbox,text width=62mm,below=of r] (m)
        {④ 证明后端（可插拔）：B1 兜底 / B2 模型采样 $k$ 篇 / B3 搜索 / B4 外部接口};
  \node[flowbox,text width=62mm,below=of m] (v)
        {⑤ 内核验证：赋值无 \texttt{sorry} 与残余元变量，交内核复核};

  \node[sidebox,right=8mm of g] (bad) {不合法\\则报错退出};
  \node[sidebox,right=8mm of v] (ok) {通过：\\⑥ 输出证明\\与验证记录};
  \node[sidebox,right=8mm of m,yshift=-14mm] (rep)
        {⑦ repair：失败码 $+$\\错误原文回灌};
  \node[sidebox,below=7mm of rep,text width=32mm] (dec)
        {⑧（可选）引理分解：\\现场造引理，\\证它后回原题};

  \draw[arr] (in)--(g);  \draw[arr] (g)--(t); \draw[arr] (t)--(r);
  \draw[arr] (r)--(m);   \draw[arr] (m)--(v);
  \draw[arr] (g.east)--(bad.west);
  \draw[arr] (v.east)--(ok.west);
  \draw[arr] (v.south) -- ++(0,-4mm) -| (rep.south);
  \draw[arr] (rep.north) |- (m.east);
  \draw[arr] (rep.south) -- (dec.north);
  \draw[arr] (t.west) -- ++(-7mm,0) |- (v.west);
  \node[left=1mm of v,font=\footnotesize,text=sggray,xshift=2mm,yshift=3mm]{命中则直达};
\end{tikzpicture}
\end{adjustbox}
\caption{在线求解流程：从命题到通过内核检验的证明。}
\end{figure}

\begin{table}[htbp]
\centering
\small
\begin{tabularx}{\textwidth}{@{}p{0.8cm}p{3.4cm}Y@{}}
\toprule
\textbf{步骤} & \textbf{模块} & \textbf{做什么与产出什么} \\
\midrule
① & 门检（Lean） & 判断命题是否为合法命题。不合法直接报错退出，
这是题目的问题，而不是“证不出来” \\
\midrule
② & 廉价 tactic 兜底（Lean） & 在预算内依次尝试若干自动化 tactic。
命中时该 tactic 本身即为证明（很多题根本不需要调用模型） \\
\midrule
③ & 分层检索（新增） & 第一层在引理库中按相关度取前若干条；库里没有的，
回落到 Mathlib 前提检索 \\
\midrule
④ & 证明后端（可插拔） & 统一接口：给定命题，返回若干候选证明脚本。
后端可增删替换，例如本地 tactic、语言模型采样、证明搜索、外部自动证明器 \\
\midrule
⑤ & 内核验证（Lean） & 造孤立义务目标，执行候选证明脚本，
最后检查赋值中无 \texttt{sorry} 与残余元变量，并交内核复核 \\
\midrule
⑥ & 输出 & 仅输出通过 ⑤ 的证明脚本，并附验证记录 \\
\midrule
⑦ & repair（新增） & 全部候选失败时，把结构化失败码与 Lean 错误原文回灌给证明后端重试，
受轮数与 token 预算限制 \\
\midrule
⑧ & 引理分解（可选） & 针对本题现场生成辅助引理并证明，再回到原题。
这些临时引理\textbf{不进入全局库} \\
\bottomrule
\end{tabularx}
\end{table}

\textbf{四个设计要点}：

\begin{enumerate}[leftmargin=1.6em,itemsep=3pt]
  \item \textbf{后端可插拔}。求解方法不固定，任何候选证明来源都可以接入，
        因为所有后端共享同一道内核闸门，正确性不依赖后端。
  \item \textbf{级联按成本排序}。先执行零成本的兜底，再执行有调用成本的模型采样，
        最后才考虑高成本的搜索。级联顺序由各后端的“边际正确率 / 边际成本”决定。
  \item \textbf{分层检索而非二选一}。自建库是精排层，Mathlib 是兜底层。
        库未覆盖的题型会自动回落到 Mathlib，系统不会因库的不足而失效。
  \item \textbf{失败信息是资产}。内核返回的是结构化失败码
        （未知标识符、类型错误、目标未闭合、含 \texttt{sorry}、超时等）与错误原文，
        比单纯的“失败”信息量大得多，可直接用于定向重试。
\end{enumerate}

\subsection{执行流程二：库的建立与管理}

\subsubsection{建立流程}

库不在解题时现造，而是由课程集离线跑闭环产出。

\begin{figure}[htbp]
\centering
\begin{adjustbox}{max width=\textwidth}
\begin{tikzpicture}[font=\small,node distance=5.5mm]
  \node[offbox,text width=78mm] (c) {课程集 $C$（与评测集\textbf{不同源}）};
  \node[offbox,text width=78mm,below=of c] (s1) {① 采轨迹：记录证明过程中的子目标签名};
  \node[offbox,text width=78mm,below=of s1] (s2) {② 需求挖掘：$d(g) = \text{freq}(g) \times \text{cost}(g)$，
        跨不少于 $m$ 个目标才算需求};
  \node[offbox,text width=78mm,below=of s2] (s3) {③ 引理生成：条件化在未解目标 $+$ 需求 $+$ 库中范例};
  \node[offbox,text width=78mm,below=of s3] (s4) {④ 门检与硬门：合法 且 非平凡 且 新颖};
  \node[offbox,text width=78mm,below=of s4] (s5) {⑤ 求解与验证：至少一篇候选通过内核（可证硬门）};
  \node[offbox,text width=78mm,below=of s5] (s6) {⑥ 软分：压缩收益与依赖关系测量};
  \node[offbox,text width=78mm,below=of s6] (s7) {⑦ 选择：库容约束下的覆盖最大化（子模贪心）};
  \node[offbox,text width=78mm,below=of s7] (s8) {⑧ 物化入库：生成有名常量并编译校验};

  \draw[arr] (c)--(s1); \draw[arr] (s1)--(s2); \draw[arr] (s2)--(s3);
  \draw[arr] (s3)--(s4); \draw[arr] (s4)--(s5); \draw[arr] (s5)--(s6);
  \draw[arr] (s6)--(s7); \draw[arr] (s7)--(s8);
  \draw[dasharr] (s8.east) -- ++(10mm,0) |- (s3.east)
        node[pos=0.35,right=0.5mm,font=\footnotesize,text=sggray]{下一轮};
\end{tikzpicture}
\end{adjustbox}
\caption{库的建立流程：课程集经过八步转化为可编译的引理库。}
\end{figure}

闭环中，三个角色分工明确：

\begin{itemize}[leftmargin=1.6em,itemsep=2pt]
  \item \textbf{求解者}在课程集上生成证明并留下轨迹；
  \item \textbf{出题者}针对需求缺口生成引理候选；
  \item \textbf{评审者}的职责由形式化判据承担，即硬门与软分，决定哪些候选入库。
\end{itemize}

注意：出题者与评审者\textbf{只存在于离线环节}；在线解一道题时它们不出场。

\subsubsection{库的管理规则}

\textbf{（一）两种来源}。库同时接收两类引理：\emph{精选层}，即从 Mathlib 检索筛选出
对特定场景特别有用的引理，成本低、优先做；\emph{自造层}，即过程中新证明出来的命题，
成本高、后做，但它是产生新数学内容的唯一入口。

\textbf{（二）四条准入规则}：

\begin{enumerate}[leftmargin=1.6em,itemsep=2pt]
  \item 自造引理必须通过同一套硬门。“已经被证明”只满足可证这一条，
        还要通过非平凡与新颖的检验。
  \item 必须记录来源信息（在哪道题、哪一轮产生），否则无法在评测时判断是否发生泄漏。
  \item “常用”必须有实测依据，即实测被依赖率与覆盖增益，不能凭印象判断。
  \item 入库前先与 Mathlib 去重。若自造引理与 Mathlib 中已有命题等价，
        不必重复入库，而应把这条对应关系存为\textbf{推荐映射}，
        告诉求解器“这类题优先考虑那条引理”。这本身比存一条重复引理更有价值。
\end{enumerate}

\textbf{（三）权重与排序}。“常用”不能直接等同于被引用次数：若权重只按引用次数累计，
而引用机会又取决于排序，就会形成正反馈，使少数几条垄断呈现机会、新引理永远上不来。
因此采用
\[
\text{权重} \;=\; \frac{\text{被引用次数}}{\text{被展示次数}}
\;+\; \text{探索额度}.
\]
即类似点击率而非单纯点击量，并给新入库引理一个初始展示额度，保证它有机会被验证是否有用。

\textbf{（四）淘汰与版本}。库不能只增不减。需要定期在工作负载上重跑评测，
把权重低于阈值的引理降级到冷存（保留记录，不直接删除）。每次入库同时保存库的版本快照，
否则实验不可复现。

\textbf{（五）六条污染路径与对策}：

\begin{table}[htbp]
\centering
\small
\begin{tabularx}{\textwidth}{@{}p{4.4cm}Y@{}}
\toprule
\textbf{污染路径} & \textbf{对策} \\
\midrule
重复入库：同一道题反复运行，等价引理反复进库 & 新颖性硬门对\textbf{库内}去重 \\
\midrule
过拟合少数题：库中引理全部围绕两三道题 & 准入门槛加跨目标计数，一条引理至少要能帮到两个不同目标 \\
\midrule
语料来源单一：课程集风格高度同质 & 课程集按子领域与题型分层抽样 \\
\midrule
过拟合工作负载：反复用同一批题挑选引理 & 工作负载只用于挑选，并留出一部分不参与挑选、只用于验证 \\
\midrule
评测污染：用测试集建库 & 数据角色硬性隔离，并记录来源信息 \\
\midrule
自我强化：呈现机会被少数引理垄断 & 权重按展示次数归一，并保留探索额度 \\
\bottomrule
\end{tabularx}
\end{table}

\textbf{（六）冻结原则}。评测时库冻结，解题过程中不再更新；
在线环节现场生成的临时引理不进入全局库。否则所测得的增益属于自我循环，结论不可用。

%======================================================================
\section{工具模块}

\subsection{模块总览}

下表中，Lean 侧模块的路径以 \lp{sgslean/} 为根，Python 侧模块的路径以 \lp{sgsr/} 为根，
实验脚本位于仓库的 \lp{scripts/} 目录下。

\begin{table}[htbp]
\centering
\small
\begin{tabularx}{\textwidth}{@{}p{2.2cm}p{5.0cm}Y@{}}
\toprule
\textbf{模块} & \textbf{位置} & \textbf{职责与状态} \\
\midrule
门检 & \lp{SgsLean/Gate.lean} & 判定命题是否合法；已实现 \\
\midrule
内核验证 & \lp{SgsLean/Verify.lean} & 整篇证明的重放与内核终检；已实现 \\
\midrule
非平凡性 & \lp{SgsLean/Trivial.lean} & 判定命题能否被自动化 tactic 秒杀；
已实现（可复用为在线兜底臂） \\
\midrule
新颖性 & \lp{SgsLean/Novelty.lean} & 与父目标及库内命题的等价性查重；已实现 \\
\midrule
软分测量 & \lp{SgsLean/Measure/} & 压缩收益与依赖关系抽取；已实现，接线中 \\
\midrule
物化与库 & \lp{SgsLean/Materialize.lean}\newline\lp{pipeline/library.py} & 生成有名常量、
编译校验、库的读写与统计；已实现 \\
\midrule
轨迹与需求 & \lp{SgsLean/Trace.lean}\newline\lp{pipeline/demand.py} &
子目标签名与需求挖掘；已实现 \\
\midrule
覆盖与选择 & \lp{pipeline/coverage.py} & 覆盖度计算与子模贪心选择；已实现 \\
\midrule
闭环编排 & \lp{pipeline/runner.py} & 十步一轮的建库编排；已实现（离线冒烟通过） \\
\midrule
模型服务 & \lp{models/} & 真代理、假服务、提示词与解析、OpenAI 兼容后端；已实现 \\
\midrule
Lean 服务 & \lp{SgsLean/Server.lean}\newline\lp{verification/client.py} &
标准输入输出 JSON 协议，支持批处理与预算配置；已实现 \\
\midrule
\textbf{证明器入口} & \lp{pipeline/prover.py}\newline\lp{scripts/prove.py} &
\textbf{待建}：单题进、证明出的统一入口 \\
\midrule
\textbf{检索层} & \lbrack 位置待定 \rbrack & \textbf{待建}：Mathlib 前提检索与库的按需取用 \\
\midrule
\textbf{repair 回路} & \lp{models/prompts.py} & \textbf{待建}：失败信息的回灌与重试 \\
\bottomrule
\end{tabularx}
\end{table}

\subsection{Lean 侧工具}

Lean 侧是本项目的正确性基础，位于 \texttt{sgs-reap/sgslean/}，以 lake 包形式组织，
依赖 Mathlib。它以标准输入输出 JSON 服务的形式对外提供能力：

\begin{itemize}[leftmargin=1.6em,itemsep=2pt]
  \item 命题合法性检查与命题到类型表达式的 elaborate；
  \item 整篇证明的验证（含内核终检）；
  \item 非平凡性与新颖性判定；
  \item 压缩收益与依赖关系测量；
  \item 已验证引理的物化（生成有名常量）与编译校验；
  \item 子目标轨迹记录。
\end{itemize}

验证链路的关键在于：候选证明先被包装成 tactic 执行，随后检查赋值中是否存在
\texttt{sorry} 与残余元变量，最后交给 Lean 内核复核。需要特别指出，在本版 Lean 中
\texttt{sorry} 的语法节点与上游守卫所匹配的名称不一致，仅靠语法检查会漏判，
真正拦住它的是内核终检这一步。

\subsection{Python 侧工具}

Python 侧位于 \texttt{sgs-reap/sgsr/}，按职责分为四个子包：

\begin{itemize}[leftmargin=1.6em,itemsep=2pt]
  \item \texttt{models/}：模型服务。包含真实代理（提供生成、评审、求解等端点）、
        确定性假服务（离线测试用）、提示词与解析器、OpenAI 兼容后端、配置读取。
  \item \texttt{pipeline/}：流水线。包含需求挖掘、引理生成组织、库的读写、
        覆盖度与选择、以及闭环编排。
  \item \texttt{verification/}：Lean 服务的常驻客户端，负责批处理与预算传递。
  \item \texttt{data/} 与 \texttt{utils/}：数据模式定义，以及服务启动与日志辅助。
\end{itemize}

\subsection{复用与上游}

项目在 \texttt{sgs-reap/reap-fork/} 中维护了一个上游证明器项目的本地分支。
其中\textbf{验证内核是本项目的地基}：Lean 侧的门检、验证、轨迹与非平凡性判定
都直接建立在它的证明重放与内核检查原语之上。此外，上游完整的证明搜索实现
（含搜索树与超参数配置）已经就绪，但默认关闭，作为可选的证明后端保留。

\subsection{待建模块}

\begin{enumerate}[leftmargin=1.6em,itemsep=3pt]
  \item \textbf{证明器入口}：接收一条命题或一个 \texttt{.lean} 文件，
        串起“门检、兜底、检索、后端采样、验证、重试”，并写回通过的证明。
        这是本项目作为产品的本体，同时也是最终评测的入口。
  \item \textbf{检索层}：现状是将库整体注入提示词，这在库规模增大后不可行
        （上下文装不下，且无关引理会变成噪声）。需要改为按相关度取前若干条，
        并接入 Mathlib 前提检索作为兜底。
  \item \textbf{repair 回路}：把结构化失败码与错误原文写回提示词，驱动定向重试。
        失败信息在上游已经产出，目前尚未被利用。
  \item \textbf{兜底臂}：把非平凡性判定所用的一组 tactic 从“返回是否秒杀”
        改造为“返回证明脚本”，使其在在线流程中直接产成证明。
\end{enumerate}

\subsection{数据角色与评测}

数据按角色严格隔离，这条规则不可违反：

\begin{table}[htbp]
\centering
\small
\begin{tabularx}{\textwidth}{@{}p{2.4cm}p{3.4cm}Y@{}}
\toprule
\textbf{数据集} & \textbf{角色} & \textbf{约束} \\
\midrule
课程集 $C$ & 建库：需求挖掘、引理生成、入库 & 必须与评测集不同源；
这是唯一可以用于建库的数据 \\
\midrule
开发集 $D$ & 调参：选择工作负载、确定采样数、观察方向 & 可用于反复实验，
但报告中必须说明所有超参数均在开发集上确定 \\
\midrule
测试集 $T$ & 最终评测 & 框架与超参数冻结后只运行一次，结果无论好坏均如实报告 \\
\bottomrule
\end{tabularx}
\end{table}

评测按\textbf{领域分桶}报告正确率，以便观察在未见领域上的表现是否存在塌陷。
需要提前注意一个来源风险：常用竞赛数学基准中的部分题目来自公开数学题库，
因此该题库不能用作课程集；另一个候选语料来源与竞赛题库存在聚合关系，
使用前必须先与评测集查重。若不做这一步，所测得的增益将因数据泄漏而不可用。

关于语料的规模与难度，有一条来自实测的经验：库能起作用的地方是求解器
\textbf{差一口气}的地方。完全解不出的题目上无法测出增益，过于简单的题目又不需要引理。
真正有信号的区间是解出率介于零与一之间的“近失手”题目，而这类题在常用基准中
占比约为一成。因此课程集的筛选应当以解出率分布为依据，把难度瞄准在
求解器能力边界附近，而非一味追求难度。

\appendix
\section{术语对照}

为避免同名异义，下表给出本项目与 SGS 之间关键术语的对应关系。

\begin{table}[htbp]
\centering
\small
\begin{tabularx}{\textwidth}{@{}p{2.6cm}p{4.6cm}Y@{}}
\toprule
\textbf{术语} & \textbf{在 SGS 中} & \textbf{在本项目中} \\
\midrule
求解者 & 语言模型生成整篇证明 & 相同，仍然是语言模型。Lean 不是求解者，
而是裁判 \\
\midrule
评审者 & 语言模型对合成题打分 & 被替换：由 Lean 计算的形式化判据，
即硬门乘以软分，另含两臂覆盖测量 \\
\midrule
出题者 & 条件化在未解目标上生成训练题 & 生成引理候选，输入扩展为需求与库中范例 \\
\midrule
验证 & 由 Lean 编译器判定 & 相同，但更严格：含内核终检与元变量检查 \\
\midrule
循环粒度 & 一批训练题，通过权重更新变强 & 分为两个尺度：
离线建库循环（一批课程集，更新库）与在线求解（单题一趟加修复） \\
\midrule
反馈流向 & 评审分数进入训练，更新出题者权重 & 分数流向选择层，决定哪些引理入库；
回灌给出题者的只有库中范例与需求 \\
\bottomrule
\end{tabularx}
\end{table}

\end{document}
